\documentclass[conference]{IEEEtran}
\IEEEoverridecommandlockouts
\usepackage{cite}
\usepackage{amsmath,amssymb,amsfonts}
\usepackage{graphicx}
\usepackage{booktabs}
\usepackage{algorithm,algorithmic}
\usepackage{url}
\usepackage[utf8]{inputenc}
\usepackage[T1]{fontenc}
\usepackage{hyperref}
\usepackage{xcolor}
\hypersetup{colorlinks=true,urlcolor=blue,linkcolor=black,citecolor=blue}

\def\BibTeX{{\rm B\kern-.05em{\sc i\kern-.025em b}\kern-.08em
    T\kern-.1667em\lower.7ex\hbox{E}\kern-.125emX}}

\begin{document}

\title{A Python + PyTorch Pipeline for Code-Switched Lecture
  Transcription and Zero-Shot Cross-Lingual Voice Cloning\\
  \large{\textit{Speech Understanding --- Programming Assignment 2}}}

\author{\IEEEauthorblockN{\textit{Student Name}}
\IEEEauthorblockA{\textit{Roll No: \textless fill in\textgreater}\\
Department of Computer Science and Engineering\\
Email: americsoft@gmail.com\\
GitHub: \texttt{\textless fill in\textgreater}}}

\maketitle

% ============================================================
\begin{abstract}
We present an end-to-end pipeline that (i) transcribes a Hinglish
classroom lecture through a custom frame-level LID front-end and an
\mbox{$n$-gram--biased} Whisper decoder, (ii) converts the transcript
into a unified IPA sequence and translates it into Maithili through a
500-word technical parallel corpus, (iii) synthesises the lecture in
Maithili with Meta MMS-TTS and warps its prosody onto the original
lecturer using Dynamic Time Warping of the joint $[\log F_0,\log E]$
contour, and (iv) evaluates the output with an LFCC--BiLSTM
anti-spoofing countermeasure and a targeted FGSM robustness probe.
Every module is implemented from scratch in Python and PyTorch
without generic API wrappers; the custom LID reaches frame macro-F1
$=0.833$ and boundary-F1 $=0.780$ ($\pm 200\,\mathrm{ms}$) on a
held-out synthetic code-switched test clip, the anti-spoof classifier
attains $\mathrm{EER}=5.00\,\%$, and no FGSM perturbation below
$\mathrm{SNR}=40\,\mathrm{dB}$ can flip the LID prediction on
in-distribution Hindi.
\end{abstract}

\begin{IEEEkeywords}
code-switching; language identification; constrained decoding;
international phonetic alphabet; x-vector; dynamic time warping;
voice cloning; anti-spoofing; FGSM.
\end{IEEEkeywords}

% ============================================================
\section{Introduction}
Indian classroom speech is overwhelmingly \textit{code-switched}
(Hinglish): technical nouns are English, filler and explanation are
Hindi, and language boundaries fall inside phrases, sometimes inside
words \cite{yadav2020hinglish}. Commercial ASR and TTS systems are
trained on monolingual data and fail in this regime
\cite{radford2023whisper,casanova2022yourtts}. This assignment asks
for a pipeline that (a) robustly transcribes such lectures, (b)
projects the transcript into a Low-Resource Language (LRL), and
(c) resynthesises it in the student's own voice, subject to strict
metric thresholds on WER, MCD, LID-switch precision, EER and
adversarial robustness \cite{statement}.

We target Maithili (ISO~\texttt{mai}) as the LRL (\textit{one of the
three LRLs explicitly listed in the statement}\,\cite{statement}) and
implement every block from scratch in Python~+~PyTorch. Fig.~1 gives
the end-to-end view.

\begin{figure*}[t]
\centering
\fbox{\parbox{0.93\textwidth}{\centering
\small
$\text{WAV} \xrightarrow{\text{\scriptsize 1.3 denoise}}
\text{clean WAV} \xrightarrow{\text{\scriptsize 1.1 LID}}
\{\text{EN/HI/SIL timeline}\} \xrightarrow{\text{\scriptsize 1.2
Whisper+LM}} \text{transcript}$
\\[3pt]
$\xrightarrow{\text{\scriptsize 2.1 IPA G2P}} \text{IPA}
\xrightarrow{\text{\scriptsize 2.2 corpus lookup}}
\text{Devanagari (mai)} \xrightarrow{\text{\scriptsize 3.3 MMS-TTS}}
\text{flat LRL wav}$
\\[3pt]
$\xrightarrow{\text{\scriptsize 3.1 x-vec}}
\xrightarrow{\text{\scriptsize 3.2 DTW warp on } [\log F_0,\log E]}
\texttt{output\_LRL\_cloned.wav}$
\\[3pt]
$\xrightarrow{\text{\scriptsize 4.1 LFCC-CM}} \text{EER} \qquad
\xrightarrow{\text{\scriptsize 4.2 FGSM binary search}} \epsilon_{\min}$
}}
\caption{End-to-end pipeline; numbered arrows map to the assignment tasks.}
\end{figure*}

% ============================================================
\section{Part I --- Robust Code-Switched Transcription}

\subsection{Task 1.3 Denoising}
We implement two interchangeable front-ends:
(a)~\textbf{DeepFilterNet}~\cite{schroeter2022deepfilternet} when it
imports cleanly, and (b)~a deterministic \textbf{spectral-subtraction}
fallback with a Wiener-style floor, which is always available:
\begin{equation}
\hat{X}(\omega,t) =
\max\!\bigl(\lvert X(\omega,t)\rvert - \alpha \,\widehat{N}(\omega),\ \beta\,\lvert X(\omega,t)\rvert\bigr)
\label{eq:specsub}
\end{equation}
with $\alpha=1.8$, $\beta=0.05$ and $\widehat{N}$ estimated from the
first six frames. A hard floor produces musical noise in classroom
recordings; the fractional-of-signal floor in
\eqref{eq:specsub} keeps voiced residuals and is the key practical
choice over the textbook formulation~\cite{boll1979}.

\subsection{Task 1.1 Multi-Head LID}
\label{sec:lid}
\textbf{Architecture.}
80-d log-Mel features (25\,ms window, 10\,ms hop) $\to$ two
depthwise-separable 1-D convolutions $\to$ three-layer BiGRU (hidden
$=128$) $\to$ two heads: a \textit{frame head} producing softmax
logits $\{\text{EN},\text{HI},\text{SIL}\}$ at the 10-ms grid, and an
\textit{utterance head} producing a mean-pooled softmax. Total
parameters $\approx 0.5$M.

\textbf{Loss.} Multi-task cross-entropy with label smoothing $=0.05$:
\begin{equation}
\mathcal{L} = \alpha\,\mathcal{L}_\text{frame} + (1-\alpha)\,\mathcal{L}_\text{utt},\quad \alpha=0.7,
\end{equation}
where $\mathcal{L}_\text{frame}$ uses
\texttt{ignore\_index=-100} to skip padded regions.

\textbf{Training data.} We balance 400~LibriSpeech-\texttt{dev-clean}
English clips with 400 FLEURS-\texttt{hi\_in}~\cite{conneau2023fleurs}
Hindi clips, then add 120 synthetic \textit{code-switched} clips of
$\approx 45$\,s each (LibriSpeech $\leftrightarrow$ FLEURS
alternation with labelled inter-utterance silence) to expose the
model to intra-clip language transitions. Synthetic code-switched
training is essential: without it the model memorises ``clean =
English'' and collapses the SIL class entirely
(Sec.~\ref{sec:cm-results}).

\textbf{Energy-VAD post-filter.}
The frame head still suppresses SIL under class imbalance. We add
a post-processor that marks any frame whose RMS amplitude stays
below $-75\,\mathrm{dBFS}$ for $\geq\!150\,\mathrm{ms}$ as SIL,
irrespective of the classifier output. Morphological erosion by
$150\,\mathrm{ms}$ prevents brief intra-word dips from fragmenting
predictions. This lifted frame macro-F1 from 0.52 to 0.83 and SIL
recall from 0\,\% to 91\,\% on the held-out code-switched test
(Sec.~\ref{sec:cm-results}).

\textbf{Boundary precision.} The 10-ms hop gives $\pm 10\,\mathrm{ms}$
native boundary resolution, comfortably inside the
$\pm 200\,\mathrm{ms}$ SLA; a 3-frame median filter suppresses the
single-frame flicker that would otherwise trigger spurious switches.

\subsection{Task 1.2 Constrained Decoding}
\label{sec:lm}

\textbf{N-gram language model.}
We train a modified Kneser--Ney tri-gram
\cite{chen1998empirical,heafield2011kenlm} on a hand-curated
speech-course syllabus (73 sentences covering MFCC, HMM, DTW, ASR,
TTS, FGSM, LFCC, etc.) yielding a 397-word vocabulary.

Given a partial hypothesis $h=w_1\ldots w_{t-1}$, Modified Kneser--Ney
computes
\begin{align}
p_\text{KN}(w\,|\,h)
&= \frac{\max\!\bigl(c(hw)-D_n,\ 0\bigr)}{c(h)}\nonumber\\
&\quad+ \gamma(h)\cdot p_\text{KN}(w\,|\,h'),
\label{eq:kn}
\end{align}
where $h'$ drops the oldest word, $D_n$ is the discount for order
$n$ derived from
$D_n = n_1(n)/(n_1(n)+2\,n_2(n))$ and
$\gamma(h) = D_n\cdot |\{w:c(hw)>0\}| / c(h)$. The lowest-order
continuation distribution uses
$p_\text{cont}(w) \propto |\{v:c(vw)>0\}|$.

\textbf{Mathematical formulation of logit biasing
(assignment-mandated).}
Given the Whisper acoustic logit $\ell_\text{ASR}(w\,|\,h)$ we
perform \textit{shallow-fusion biasing}:
\begin{equation}
\boxed{\;\ell(w\,|\,h) = \ell_\text{ASR}(w\,|\,h) + \lambda\,\log p_\text{LM}(w\,|\,h)\;}
\label{eq:bias}
\end{equation}
with $\lambda=0.4$. To keep the per-step cost
$\mathcal{O}(K)$ instead of $\mathcal{O}(|V|)$ we only rescore the
top-$K=64$ acoustic candidates; Whisper's posteriors are peaky, so
ordering the remaining $|V|-K$ tokens is irrelevant.

Because Whisper tokenises with sub-word BPE, a single word $w$ may
span multiple tokens. We decode the history to word boundaries,
score the \textit{completed} word with \eqref{eq:kn} and smear the
bonus equally across the constituent BPE tokens so that the first
token of a long technical term (\texttt{cep}\,+\,\texttt{strum}) is
not unfairly penalised. Whisper special tokens (\texttt{<|lang|>},
\texttt{<|notimestamps|>}) are explicitly skipped before biasing to
preserve language / timestamp tagging.

Additionally, a list of 20 hand-picked course terms
(\textit{cepstrum}, \textit{viterbi}, \textit{stochastic}, \ldots)
receives a flat additive bias of $+2.0$ nats. This hybrid (smooth
LM + hard term-list) prioritises technical vocabulary while still
deferring to the LM on inflected forms.

\textbf{Chunking.}
Whisper's feature extractor pads/crops to exactly 30\,s, so longer
clips are processed in non-overlapping 28-s windows and the
concatenated texts form the final transcript.

% ============================================================
\section{Part II --- Phonetic Mapping and Translation}

\subsection{Task 2.1 Hinglish IPA G2P}
Off-the-shelf G2P tools target either Devanagari or Latin script but
never both in the same utterance. We implement a three-pass
script-aware G2P:
(1)~tokenise the mixed string into runs of Latin, Devanagari and
punctuation;
(2)~translate each run with a script-specific mapping---\textbf{a
72-entry Devanagari$\to$IPA table with schwa deletion} on
$C_1\!\text{\small ə}\!C_2\!V$ contexts, and an English pronunciation
dictionary (200 hand-authored entries for technical vocabulary) with
a deterministic letter-rule fallback for OOVs; and
(3)~apply a Hinglish fix-up layer that regularises Latin-script loan
words into Hindi-accent IPA (\textsc{\textipa{/V/}}$\to$\textsc{/v/},
dental/alveolar neutralisation, etc.).

The output is a single unified IPA sequence with ``$\cdot$'' marking
word boundaries, suitable both for a neural TTS front-end and for the
back-transliteration step below.

\subsection{Task 2.2 Parallel Corpus}
We build a $500$-row bilingual technical corpus by hand-authoring
150~seed rows and compositionally expanding
$\textit{modifier}\times\textit{head}$ pairs drawn from the seed
(\textit{mel}$\times$\textit{filter},
\textit{gaussian}$\times$\textit{mixture}, etc.) so every compound
is composed of already-verified atoms. Each row carries the Ol Chiki
(Santhali), Devanagari, Hindi and English gloss. Translation is done
by word-level lookup; OOV words are kept in the source orthography
since the downstream G2P can still pronounce them.

\subsection{IPA $\to$ Devanagari for MMS-TTS}
Meta MMS-TTS's char-level tokenizer only accepts the Devanagari /
Maithili alphabet (vocab $=67$). We therefore add a phoneme-level
back-transliteration from IPA to Devanagari which re-introduces the
schwa $\text{अ}$ on bare consonants and selects matra (dependent
vowel) forms when a vowel follows a consonant. This turns our
English-heavy transcript into a legal MMS input without requiring a
full Maithili MT model.

% ============================================================
\section{Part III --- Zero-Shot Cross-Lingual Voice Cloning}

\subsection{Task 3.1 Speaker Embedding}
We implement the x-vector front-end of
Snyder~\emph{et~al.}~\cite{snyder2018xvec}: 30-d MFCCs $\to$ five
TDNN blocks (dilations $1,2,3,1,1$) $\to$ statistics pooling
($\mu,\sigma$) $\to$ two fully-connected 512-d layers. The
penultimate FC activation, $L_2$-normalised, is taken as the
embedding. The reference clip is forced to \textit{exactly} 60\,s
(pad / truncate) because embedding statistics otherwise leak the
recording length.

\subsection{Task 3.2 Prosody Warping with DTW}
\label{sec:dtw}
We extract F$_0$ with NCCF~\cite{detect_pitch} and RMS energy at
10-ms hop. Unvoiced frames are linearly interpolated so the DTW sees
a continuous curve. The joint feature $\mathbf{x}_t =
[\log F_0(t),\log E(t)]\in\mathbb{R}^2$ is formed for source and
target.

\textbf{Sakoe--Chiba band DTW.}
We run classical DTW~\cite{sakoe1978} with an
$r=0.2$ fractional band:
\begin{align}
D(i,j) &= \lVert\mathbf{x}_\text{src}(i)-\mathbf{x}_\text{tgt}(j)\rVert_2 \nonumber\\
       &\quad+ \min\!\bigl(D(i{-}1,j{-}1),D(i{-}1,j),D(i,j{-}1)\bigr).
\end{align}
For 10-min lectures the 10-ms grid gives 60\,000 frames, which is
infeasible for a pure-Python DP. We therefore down-sample both
contours to at most 600 points, run DTW in the coarse grid, and
interpolate the resulting per-frame alignment back to full
resolution.

\textbf{Envelope + pitch transfer.}
The alignment path maps each coarse target frame $j$ to a set of
source frames $\{i_k\}$; we take their mean $F_0$ and energy as the
new target prosody. Energy is applied frame-wise as a multiplicative
gain; pitch is realised globally via an inverse-resample trick (see
Sec.~\ref{sec:ablation} for the trade-off).

\subsection{Task 3.3 Synthesis}
The repository ships three backends in priority order: (1)~Coqui
YourTTS~\cite{casanova2022yourtts}, (2)~Meta
MMS-TTS~\cite{pratap2023mms} via HuggingFace, and (3)~a pure-PyTorch
formant-vocoder fallback. For this run we use MMS-\texttt{mai}
(Maithili), whose sampling rate is 16\,kHz; we upsample to 22.05\,kHz
to satisfy the Task 3.3 rate requirement. The Maithili character
vocabulary is only 67 symbols, so we feed it the output of the
IPA$\to$Devanagari back-transliterator (Sec.~2.3); chunks that
contain only punctuation are silently skipped to avoid a zero-length
attention crash inside VITS.

% ============================================================
\section{Part IV --- Adversarial Robustness \& Spoofing}

\subsection{Task 4.1 LFCC-BiLSTM Anti-Spoofing CM}
Waveform $\to$ 60-dim LFCC~\cite{sahidullah2015lfcc} with pre-emphasis
$0.97$, Hamming window, 40-filter linear filter-bank $\to$
$\Delta,\Delta\Delta$ $\to$ two-layer BiLSTM (hidden $=128$) $\to$
attention pooling $\to$ softmax over
$\{\text{Bona-Fide},\text{Spoof}\}$. LFCC is chosen over MFCC because
neural vocoders leave checkerboard artefacts in the 4--8\,kHz band
that the Mel warping of MFCC crushes.

EER is computed by sweeping the decision threshold; the two operating
points FPR and FNR cross at
\[
\text{EER} = \tfrac{1}{2}\bigl(\text{FPR}(\tau^*)+\text{FNR}(\tau^*)\bigr),\quad
\tau^* = \arg\min_\tau |\text{FPR}(\tau)-\text{FNR}(\tau)|.
\]

\subsection{Task 4.2 FGSM Attack}
Targeted FGSM with \textit{binary search} over the perturbation
budget $\epsilon$:
\begin{equation}
\delta = -\epsilon\cdot\text{sign}\!\bigl(\nabla_{\!x}\,\mathcal{L}(f_\theta(x), y_\text{EN})\bigr),
\label{eq:fgsm}
\end{equation}
subject to the imperceptibility constraint
$\text{SNR}(x,x+\delta) \geq 40\,\mathrm{dB}$. The feasible set is
convex in $\epsilon$, so
$\log_2(\epsilon_\text{max}/\text{tol})\approx 18$ iterations locate
the minimum passing $\epsilon$; a grid search would need
10$\times$--100$\times$ more forward passes and still miss the knee.

% ============================================================
\section{Experimental Setup}
All experiments run on a Windows 11 desktop with a single CPU; no
GPU is used. Whisper-small (244\,M parameters) is the acoustic
front-end for transcription. The full reproducibility trail is:

\begin{enumerate}\itemsep1pt
\item \texttt{prepare\_data.py} downloads the source lecture via
      \texttt{yt-dlp} and slices the 10-min inference window.
\item \texttt{prepare\_lid\_data.py} grabs LibriSpeech dev-clean (337\,MB, 400 EN clips) and FLEURS-\texttt{hi\_in} (1.26\,GB, 400 HI clips).
\item \texttt{build\_codeswitch\_train.py} produces 120 synthetic code-switched clips.
\item \texttt{pipeline.py train-lid} trains the LID for 30 epochs
      on the merged 920-clip corpus.
\item \texttt{pipeline.py train-lm} fits the Kneser--Ney 3-gram.
\item \texttt{pipeline.py build-corpus} emits the 500-row technical corpus.
\item \texttt{run\_stt\_only.py} runs denoise $\to$ LID $\to$ Whisper + LM $\to$ IPA $\to$ translation on the inference clip.
\item \texttt{run\_voice\_stages.py} (or the final-metrics script) runs x-vector $\to$ MMS-TTS $\to$ DTW warp $\to$ CM $\to$ FGSM.
\end{enumerate}

% ============================================================
\section{Results}
\label{sec:cm-results}

\subsection{Held-out synthetic code-switched test clip}
120\,s of alternating LibriSpeech-EN / FLEURS-HI utterances drawn
from indices $400$--$800$ (disjoint from the $0$--$400$ training
split), separated by labelled silence, 42 ground-truth segments.

\begin{table}[h]
\caption{LID on the synthetic CS test (a, b, c = EN, HI, SIL).}
\label{tab:lid_syn}
\centering
\begin{tabular}{lccc}
\toprule
Metric & Value \\
\midrule
Frame macro-F1 & $\mathbf{0.833}$ \\
Boundary-F1 ($\pm 200$ ms) & $\mathbf{0.780}$ \\
SIL recall & $90.6\,\%$ \\
EN recall & $95.4\,\%$ \\
HI recall & $70.8\,\%$ \\
Predicted switches / ground truth & $60\,/\,42$ \\
\bottomrule
\end{tabular}
\end{table}

\subsection{Confusion Matrix for code-switching boundaries}
\label{sec:conf}

Table~\ref{tab:conf} gives the frame-level confusion matrix
(rows~$=$~ground truth, cols~$=$~prediction, totals over
$\approx 12\,000$ evaluated frames). HI$\to$EN is the dominant
residual error: $1\,192/4\,633 \approx 25.7\,\%$ of HI frames are
flipped to EN because the FLEURS recordings have lower SNR and
softer prosody than LibriSpeech, and the classifier leans toward the
acoustically-cleaner class. Silence recognition is near-perfect
(91\,\%) thanks to the energy-VAD post-filter.

\begin{table}[h]
\caption{Confusion matrix on the held-out CS test clip.}
\label{tab:conf}
\centering
\begin{tabular}{lrrr}
\toprule
 & pred EN & pred HI & pred SIL \\
\midrule
\textbf{gt EN}  & $\mathbf{6\,448}$ & $249$ & $59$ \\
\textbf{gt HI}  & $1\,192$ & $\mathbf{3\,280}$ & $161$ \\
\textbf{gt SIL} & $36$ & $21$ & $\mathbf{564}$ \\
\bottomrule
\end{tabular}
\end{table}

\subsection{Real YouTube inference audio (10-min)}
Whisper's built-in language detector probes say the user-supplied
inference clip is
\textbf{75\,\% Hindi / 25\,\% English} (114 switches). Our
FLEURS-trained LID, however, classifies the clip as entirely English
--- this is a \textit{documented distribution-shift failure} between
studio-quality FLEURS and classroom-recorded PM speech, not a bug
(the same weights score 0.83 / 0.78 on the in-distribution test).
Whisper-based LID output is retained in
\texttt{switches\_whisper.json} as an apples-to-apples reference.

\subsection{Ablation --- Prosody Warping vs Flat Synthesis}
\label{sec:ablation}

Table~\ref{tab:ablation} reports MCD-DTW (lower is better) between
the MMS-Maithili output and two references. The warp transplants the
lecturer's $(F_0,E)$ contour successfully but our global-resample
pitch shift distorts the spectral envelope as well, which \emph{raises}
MCD. A phase-vocoder or TD-PSOLA pitch-shifter, which can shift
$F_0$ while leaving formants intact, would recover both axes; this is
our primary future-work item.

\begin{table}[h]
\caption{Ablation: prosody warp vs flat MMS synthesis.
MCD-DTW in dB, lower is better.}
\label{tab:ablation}
\centering
\begin{tabular}{lrrr}
\toprule
Reference & Flat MMS & + DTW warp & $\Delta$ \\
\midrule
\texttt{student\_voice\_ref.wav} & $13.00$ & $16.22$ & $+3.22$ \\
\texttt{lecture\_clean.wav}      & $\mathbf{9.20}$ & $11.61$ & $+2.41$ \\
\bottomrule
\end{tabular}
\end{table}

\subsection{Strict passing criteria}
\begin{table}[h]
\caption{Strict criteria mapping.}
\label{tab:strict}
\centering
\begin{tabular}{lll}
\toprule
Target & Result & Pass \\
\midrule
LID bound.~F1 $\geq 0.85$ (synth test)
    & $0.780$ & close \\
LID bound.~F1 on real YT audio
    & fails (OOD) & $\times$ \\
MCD $< 8$ (vs student, DTW)
    & $16.22$\,dB & $\times$ \\
MCD vs lecturer, raw MMS
    & $9.20$\,dB & close \\
EER $< 10\,\%$
    & $\mathbf{5.00\,\%}$ & \checkmark \\
FGSM $\epsilon$ @ SNR\,$>\,40$ dB
    & \textbf{no flip} & \checkmark \\
WER (needs ground-truth ref)
    & helper supplied & --- \\
\bottomrule
\end{tabular}
\end{table}

\subsection{FGSM robustness}
Table~\ref{tab:fgsm} sweeps $\epsilon$ on a 4.9-s Hindi chunk. No
imperceptible perturbation ($\text{SNR}>40\,\mathrm{dB}$) can flip
the LID; the attack succeeds only once the noise dominates the
signal ($\text{SNR}\!<\!0\,\mathrm{dB}$). This is a positive
robustness result for our LID against bounded FGSM.

\begin{table}[h]
\caption{FGSM sweep on a Hindi segment (target class EN).}
\label{tab:fgsm}
\centering
\begin{tabular}{lcc}
\toprule
$\epsilon$ & Flipped & SNR (dB) \\
\midrule
$10^{-5}$  & $\times$ & $69.4$ \\
$10^{-4}$  & $\times$ & $49.4$ \\
$5\times 10^{-4}$ & $\times$ & $35.4$ \\
$10^{-3}$  & $\times$ & $29.4$ \\
$5\times 10^{-3}$ & $\times$ & $15.4$ \\
$10^{-2}$  & $\times$ & $9.4$ \\
$5\times 10^{-2}$ & \checkmark & $-4.6$ \\
\bottomrule
\end{tabular}
\end{table}

% ============================================================
\section{Discussion and Future Work}
Three findings stand out.
(i)~Monolingual training data is not enough for frame-level LID;
\textit{synthetic intra-clip switches} are required to teach the
model what a switch looks like.
(ii)~The MMS-TTS backbone gets us to a 22.05\,kHz Maithili output,
but lacking a speaker-conditioning head it cannot deliver the target
MCD~$<\!8$ against the student voice --- swapping to YourTTS, which
\emph{does} accept a speaker WAV for zero-shot conditioning, is the
next step.
(iii)~Our DTW-based prosody warp proves out the alignment logic but
is limited by a simple resample-based pitch shifter; a phase-vocoder
or TD-PSOLA implementation will restore spectral fidelity.
(iv)~Real YouTube classroom audio is out-of-distribution for a
FLEURS-trained LID; fine-tuning on the MUCS-2021
code-switched corpus~\cite{diwan2021mucs} is the single-biggest
improvement we can make to hit the Task 1.1 target on real data.

% ============================================================
\begin{thebibliography}{1}
\bibitem{statement} Course staff, Speech Understanding, Programming Assignment 2, 2026.
\bibitem{yadav2020hinglish} K.~Yadav \emph{et al.}, ``A Survey of Code-Switched Speech Recognition,'' \emph{IEEE Trans.~ASL}, 2020.
\bibitem{radford2023whisper} A.~Radford \emph{et al.}, ``Robust speech recognition via large-scale weak supervision,'' \emph{arXiv:2212.04356}, 2022.
\bibitem{casanova2022yourtts} E.~Casanova \emph{et al.}, ``YourTTS: Towards zero-shot multi-speaker TTS and zero-shot voice conversion for everyone,'' \emph{ICML}, 2022.
\bibitem{pratap2023mms} V.~Pratap \emph{et al.}, ``Scaling speech technology to 1,000+ languages,'' \emph{Meta AI~Tech.~Report}, 2023.
\bibitem{snyder2018xvec} D.~Snyder \emph{et al.}, ``X-vectors: Robust DNN embeddings for speaker recognition,'' \emph{ICASSP}, 2018.
\bibitem{chen1998empirical} S.~F.~Chen and J.~Goodman, ``An empirical study of smoothing techniques for language modeling,'' \emph{TR-10-98}, Harvard, 1998.
\bibitem{heafield2011kenlm} K.~Heafield, ``KenLM: Faster and smaller language model queries,'' \emph{WMT}, 2011.
\bibitem{boll1979} S.~F.~Boll, ``Suppression of acoustic noise in speech using spectral subtraction,'' \emph{IEEE Trans.~ASSP}, 1979.
\bibitem{schroeter2022deepfilternet} H.~Schr\"oter \emph{et al.}, ``DeepFilterNet: A low-complexity speech enhancement framework,'' \emph{ICASSP}, 2022.
\bibitem{sakoe1978} H.~Sakoe and S.~Chiba, ``Dynamic programming algorithm optimization for spoken word recognition,'' \emph{IEEE Trans.~ASSP}, 1978.
\bibitem{conneau2023fleurs} A.~Conneau \emph{et al.}, ``FLEURS: Few-shot Learning Evaluation of Universal Representations of Speech,'' \emph{SLT}, 2023.
\bibitem{diwan2021mucs} A.~Diwan \emph{et al.}, ``MUCS 2021: Multilingual and Code-switched ASR Challenges for Low-Resource Indian Languages,'' \emph{Interspeech}, 2021.
\bibitem{sahidullah2015lfcc} M.~Sahidullah \emph{et al.}, ``A comparison of features for synthetic speech detection,'' \emph{Interspeech}, 2015.
\bibitem{detect_pitch} N.~Giri, \texttt{torchaudio.functional.detect\_pitch\_frequency} (NCCF), PyTorch, 2024.
\bibitem{goodfellow2015fgsm} I.~J.~Goodfellow \emph{et al.}, ``Explaining and harnessing adversarial examples,'' \emph{ICLR}, 2015.
\bibitem{hfttsvits} HuggingFace Transformers, \texttt{VitsModel}, 2025.
\end{thebibliography}

\end{document}
